%! Author = sriadityavedantam
%! Date = 3/25/26

\lesson{47}{Self-Study Three}{Day 47}

\begin{proposition}{
    If $f:X\to Y$ is continuous and $X$ is compact, then $f$ is uniformly continuous.
}
    Let $\eps > 0$.

    For each $c\in X$, by continuity of $f$ at $c$, there exists $\delta_c > 0$ such that
    \[
        d_X(x,c) < \delta_c \implies d_Y(f(x),f(c)) < \frac{\eps}{2}.
    \]

    Then $\{B(c,\delta_c)\}_{c\in X}$ is an open cover of $X$.

    By the Lebesgue Big Number Lemma, there exists $\delta > 0$ such that for all $x\in X$, there exists $c\in X$ with
    \[
        B(x,\delta)\subseteq B(c,\delta_c).
    \]

    Now suppose $d_X(x,y) < \delta$.

    Then $y\in B(x,\delta)\subseteq B(c,\delta_c)$, and also $x\in B(c,\delta_c)$.

    So
    \[
        d_Y(f(x),f(c)) < \frac{\eps}{2}
        \qquad\text{and}\qquad
        d_Y(f(y),f(c)) < \frac{\eps}{2}.
    \]

    Hence
    \begin{align*}
        d_Y(f(x),f(y))
        &\leq d_Y(f(x),f(c)) + d_Y(f(c),f(y))\\
        &< \frac{\eps}{2} + \frac{\eps}{2}\\
        &= \eps.
    \end{align*}

    So $f$ is uniformly continuous.
\end{proposition}

\begin{definition}[Contraction]
    A function $f:X\to X$ is a contraction if there exists $K$ with $0\leq K < 1$ such that for all $x,y\in X$,
    \[
        d(f(x),f(y)) \leq K\,d(x,y).
    \]
\end{definition}

\begin{definition}[Fixed Point]
    Given $f:X\to X$, a point $x\in X$ is a fixed point if
    \[
        f(x) = x.
    \]
\end{definition}

\begin{theorem}[Banach Fixed Point Theorem]{
    Let $(X,d)$ be a complete metric space.

    If $f:X\to X$ is a contraction, then there exists a unique fixed point.
}
    Pick $x_0\in X$.

    Define a sequence by
    \[
        x_{n+1} \coloneqq f(x_n).
    \]

    Then
    \[
        d(x_{n+1},x_n) \leq K^n d(x_1,x_0).
    \]

    For $m>n$, we have
    \begin{align*}
        d(x_m,x_n)
        &\leq \sum_{i=n}^{m-1} d(x_{i+1},x_i)\\
        &\leq \sum_{i=n}^{m-1} K^i d(x_1,x_0)\\
        &= d(x_1,x_0)\sum_{i=n}^{m-1} K^i\\
        &\leq \frac{K^n}{1-K} d(x_1,x_0).
    \end{align*}

    So $(x_n)$ is Cauchy.

    Since $X$ is complete, there exists $x\in X$ such that $x_n\to x$.

    Since $f$ is continuous, we have
    \[
        x = \lim_{n\to\infty} x_{n+1}
        = \lim_{n\to\infty} f(x_n)
        = f\left(\lim_{n\to\infty} x_n\right)
        = f(x).
    \]

    So $x$ is a fixed point.

    To prove uniqueness, suppose $y\in X$ is another fixed point.

    Then
    \begin{align*}
        d(x,y)
        &= d(f(x),f(y))\\
        &\leq K\,d(x,y).
    \end{align*}

    Since $0\leq K<1$, this implies
    \[
        (1-K)d(x,y)\leq 0,
    \]
    so
    \[
        d(x,y)=0.
    \]

    Hence $x=y$.
\end{theorem}